\documentclass[letterpaper, 10pt,colorlinks,linkcolor=cyan]{moderncv}
  \usepackage{lmodern}
  \usepackage[english]{babel}
  \usepackage{courier}
  \hyphenpenalty=100000
  \hbadness=99999
%% Themes
  \moderncvstyle{casual}
  \moderncvcolor{red}
  \nopagenumbers{}

% --- Name in red (color1) while preserving casual's font weight split ---
  \colorlet{lastnamecolor}{color1}
  \colorlet{firstnamecolor}{color1}
  % Force section titles to red
  \colorlet{sectioncolor}{color1}
  \colorlet{subsectioncolor}{color1}
  % Force horizontal section bar to red
  \colorlet{bodyrulecolor}{color1}
  % Override section style to remove small caps and use upshape
  \renewcommand*{\sectionstyle}[1]{{\Large\bfseries\upshape\textcolor{sectioncolor}{#1}}}
  % Override subsection style to remove small caps
  \renewcommand*{\subsectionstyle}[1]{{\large\bfseries\upshape\textcolor{subsectioncolor}{#1}}}
  % Reduce body font to 9pt
  \renewcommand{\normalsize}{\fontsize{9}{11}\selectfont}
%% Encoding
  \usepackage[utf8]{inputenc}
  \usepackage[T1]{fontenc}
%% Margins
  \usepackage[left=0.65in,right=0.65in,top=0.5in,bottom=1.0in]{geometry}
  \setlength{\hintscolumnwidth}{2.25cm}
%% Data
  \firstname{Daniel}
  \familyname{Hardesty Lewis}
  \address{\textsc{500~Riverside~Dr}}{\textsc{New~York,~NY~10027}}{\textsc{United~States}}
  \phone[mobile]{+1~(713)~371-7875}
  \email{daniel@summitgeospatial.com}
  \homepage{dlewis.ai}
  \social[linkedin]{dhardestylewis}
  \social[github]{dhardestylewis}

  \extrainfo{\href{https://scholar.google.com/citations?user=agBdssgAAAAJ}{Google Scholar}}
  \title{\large AI Consultant -- ML systems, HPC, \& scientific computing}
%% Define \cvdoublecolumn
  \newcommand{\cvdoublecolumn}[2]{
    \cvitem[0.75em]{}{
      \begin{minipage}[t]{\listdoubleitemcolumnwidth}#1\end{minipage}
      \hfill
      \begin{minipage}[t]{\listdoubleitemcolumnwidth}#2\end{minipage}
    }
  }

\begin{document}
  \definecolor{links}{HTML}{0645AD}
  \hypersetup{urlcolor=links}

%%%-----     cover letter     -----
  \recipient{Prof.\ Frederic T.\ Chong}{Department of Computer Science\\University of Chicago\\5730 S.\ Ellis Ave\\Chicago, IL 60637}
  \date{February 20, 2026}
  \opening{Dear Prof.\ Chong,}
  \closing{Sincerely,}
  \enclosure[Attached]{curriculum vit\ae{}}
  \makelettertitle

I am writing to propose a consulting engagement with your research group, where I can contribute hands-on AI and ML systems engineering to accelerate your work at the intersection of quantum computing and practical AI infrastructure.

Your group's focus on bridging quantum and classical computing architectures, from the EPiQC Expedition through Infleqtion's quantum software stack, aligns directly with my experience building and operating ML systems on leadership-class computing infrastructure. At the \href{https://www.tacc.utexas.edu/}{Texas Advanced Computing Center}, I spent five years orchestrating distributed workloads across 800,000-node supercomputer clusters, managing multi-million-dollar compute budgets, and co-teaching graduate courses in ML, CUDA, and C++ -- the exact stack that matters when pushing novel computing paradigms toward practical scale.

Since then, I have shipped production ML systems end-to-end at \href{https://www.homecastr.com/}{Homecastr} (deep generative models on noisy tabular data) and conducted research at Columbia on latent factor models with \href{https://alihrsa.github.io/}{Dir.\ Financial Engineering Ali Hirsa} and on DeepSeek's GRPO algorithm with \href{https://kostic.ece.columbia.edu/}{Dir.\ Electrical Engineering Zoran Kostic}. Every engagement has required bridging the gap between research ideas and reliable, scalable systems -- the same challenge your group faces in making quantum-classical hybrid architectures practical.

I would welcome the opportunity to discuss how I can contribute to your group's current projects, whether in ML infrastructure, compiler-level optimization, or systems engineering for hybrid quantum-classical workloads.

  \makeletterclosing
  \clearpage

  \vspace*{-2.5em}
  \makecvtitle

\vspace{-0.6em}
\noindent\begin{minipage}{\textwidth}
\small\raggedright
AI consultant specializing in ML systems engineering, HPC infrastructure, and bridging research to production. Hands-on experience shipping deep generative models, orchestrating distributed workloads at supercomputer scale, and teaching ML, CUDA, and C++ at the graduate level.
\end{minipage}
\vspace{0.6em}

  \section{Core Expertise}
  \cvitem{8 years}{\textbf{ML Systems Engineering}: Build and operate ML training and inference pipelines, including deep generative models, rigorous evaluation, and automated retraining}
  \cvitem{8~years}{\textbf{HPC \& Distributed Computing}: Scale batch workloads across leadership-class clusters, profile bottlenecks, and optimize latency \& throughput with explicit cost-performance tradeoffs}
  \cvitem{8~years}{\textbf{Backend \& Data}: Build schemas and ETL; ship API-backed services using Python, SQL, and Postgres/PostGIS}
  \cvitem{10~years}{\textbf{Scientific Computing}: Spatiotemporal modeling, numerical methods, and simulation-based decision-support across geospatial, financial, and climate domains}
  \cvitem{8~years}{\textbf{Technical Leadership}: Drive cross-functional execution via specifications and governance; mentor engineers and researchers through design reviews, code reviews, and instruction}

  \section{Technologies}
  \cvitem{Modeling}{PyTorch; TensorFlow; scikit-learn; Transformers; LoRA/QLoRA/PEFT}
  \cvitem{LLM apps}{embeddings; RAG; LangChain; LangGraph; DSPy; FAISS}
  \cvitem{Data stack}{Python; SQL; Pandas; Polars; PyArrow; Parquet; Postgres; PostGIS; Redis; Elasticsearch; Supabase}
  \cvitem{Infrastructure}{Linux/Unix; Bash; Git; Docker; Kubernetes; Ray; Spark; Dask; Airflow; CI/CD; observability}
  \cvitem{Systems}{CUDA; C++; MPI; OpenMP; SLURM; containerized HPC}

  
  \section{\textsc{Experience}}
  \subsection{\textsc{Professional}}
    \cventry{Apr.~2025 -- Present}{Co-founder}{\href{https://www.homecastr.com/}{Homecastr}}{New York, NY}{}{
    Consumer real-estate pricing and forecasting platform with interactive uncertainty fancharts and temporal leakage-safe expanding-window evaluation; owned end-to-end workflow from source data to web delivery.
    \begin{itemize}
        \item Deployed an NYC semi-supervised VAE current-price model using \textbf{tax and sales records} to handle \textbf{sparsity and noise}; achieved 12\% holdout error (vs.\ Zillow's 8.4\% internal apples-to-apples).
        \item Shipped a Houston diffusion-based forecasting system with fan charts; achieved as low as $\sim$8\% annualized compounding error, balancing design choices of error vs.\ trajectory realism, delivered via production web UI: \href{https://www.homecastr.com/}{homecastr.com}
        \item Refined product positioning and execution milestones via strategic feedback from a16z, MetaProp, and others.
    \end{itemize}}

    \cventry{Nov.~2024 -- Present}{Co-founder}{\href{https://www.polibom.com/}{PoliBOM}}{Seattle, WA}{}{
    AI tariff mitigation and trade compliance product for manufacturers, centered on BOM analysis, retrieval, and policy scenario simulation
    \begin{itemize}
      \item Specified an agentic workflow and schema definitions for BOM parsing, retrieval, and tariff simulation using Airflow, Elasticsearch and embeddings, and FastAPI serving backed by Supabase and Redis.
      \item Directed development of a \textbf{conversational AI interface} simplifying global trade workflows; enabled decision-makers to instantly \textbf{retrieve BOMs and receive proactive compliance alerts}.
      \item Executing customer discovery with the founder of ScoutyAI and tariff economists at Columbia Business School to validate product-market fit and pricing assumptions.
      \item Secured market validation by presenting the prototype to KPMG, pre-dating their subsequent launch of a \href{https://kpmg.com/xx/en/home/insights/2024/05/tariff-modeler.html}{commercial tariff-modeling product}.
    \end{itemize}}

    \cventry{Nov.~2023 -- Present}{Founder}{\href{https://summitgeospatial.com}{Summit Geospatial}}{New York, NY}{}{
    High-resolution elevation data and web delivery for engineering and hazard workflows.
    \begin{itemize}
      \item Engineered a statewide seamless Texas DEM mosaic, resampling 0.5~m LiDAR sources to 1.2~m via nearest-neighbor across 70+ state, federal, and local elevation datasets.
      \item Developed the web distribution platform end-to-end including data pipeline, tiling, hosting, delivery UX to support engineering and planning use cases.
      \item Commercial discussions to license elevation datasets to research and AI organizations.
    \end{itemize}}

    \cventry{Aug.~2021 -- Aug.~2023}{Senior Data Scientist \& Technical Lead}{\href{https://www.tacc.utexas.edu/}{Texas Advanced Computing Center (TACC)}}{Austin, TX}{}{
    \begin{itemize}
        \item Co-founded \href{https://www.summitgeospatial.com}{summitgeospatial.com} to commercialize and sustain state-level terrain models developed for the \$40M \href{https://tdis.io/}{Texas Disaster Information System (TDIS)}, delivering precision decision-support for civil engineering.
        \item Developed inter-agency RFP and MOU frameworks with state and federal partners, translating technical requirements into scope and milestones facilitating \href{https://storymaps.arcgis.com/stories/6227f24d6c5047f883d5c3497241f299}{Real-Time Inundation Mapping} and TDIS.
        \item Scaled climate and flood models on supercomputers, orchestrating 800,000-node distributed jobs across containerized HPC clusters while managing multi-million-dollar compute budgets.
        \item Developed efficient methods to produce high-resolution 1m flood maps from National Water Model outputs for Texas Emergency Management's Real-Time Inundation Mapping.
        \item Partnered with the US Army Corps of Engineers to statistically model compound flood hazards in coastal Texas, integrating high-resolution topographic data with hydrological models.
        \item Contributed fundamental research supporting Paola Passalacqua's \href{https://www.egu.eu/awards-medals/ralph-alger-bagnold/2022/paola-passalacqua/}{2022 EGU Bagnold Medal}, one of geomorphology's highest honors.
        % Formal title: Technical Reviewer
        % NSF Frontera leadership-class computing project (NSF award #1818253): https://www.nsf.gov/awardsearch/show-award/?AWD_ID=1818253
        % Awardees list: https://frontera-portal.tacc.utexas.edu/allocations/pathways-awardees/
        \item Served as Technical Reviewer for \href{https://frontera-portal.tacc.utexas.edu/fellowship/}{Frontera Doctoral Fellows} (2023) and \href{https://frontera-portal.tacc.utexas.edu/allocations/pathways-awardees/}{Large-Scale Computing Pathways} (2022) applications, evaluating \$10M+ in computing resource allocations on one of the world's largest supercomputers.
    \end{itemize}}
    
    \cventry{Feb.~2018 -- Jul.~2021}{Data Scientist \& Research Engineer}{\href{https://www.tacc.utexas.edu/}{Texas Advanced Computing Center (TACC)}}{Austin, TX}{}{
    \begin{itemize}
        \item Competed in \href{https://www.darpa.mil/program/world-modelers}{DARPA World Modelers} to improve disaster resiliency and food security modeling in East Africa.
        \item Mentored Petrobras engineers in deep learning and HPC to enable institutional technology transfer.
        \item Received research recognition from AAAI President Yolanda Gil during her \href{https://youtu.be/3Spk_UOHNxA?t=2377}{farewell address} for contributions to automated scientific modeling.
    \end{itemize}}


  \subsection{\textsc{Research}}

    \cventry{Jan.~2025 -- Oct.~2025}{Research assistant to Dir. Financial Engineering Ali Hirsa (Industry-sponsored research)}{Columbia University}{New York, NY}{}{
    \begin{itemize}
        \item Developed a multi-asset CVAE latent factor model with Skew-T Mixture priors; achieved 85\% \(R^2\) vs.\ commercial SaaS (75\%) and Fama-French (58\%) on backtested holdout.
        \item Derived a tractable, axiom-compliant per-feature SHAP algorithm using Hessian-based recursive clustering to provide stable temporal regime-dependent attributions for deep learning models.
    \end{itemize}}

    \cventry{May~2025 -- Aug.~2025}{Research assistant to Dir. Electrical Engineering Zoran Kostic}{Columbia University}{New York, NY}{}{
    \begin{itemize}
        \item Math benchmarking and finetuning of Deepseek's GRPO algorithm
    \end{itemize}}
    
    \cventry{July~2019 -- Aug.~2019}{Summer research intern to Pres. AAAI Yolanda Gil}{University of Southern California}{Los Angeles, CA}{}{
    \begin{itemize}
        \item Integrated hydrological models into the \href{http://mint-project.info/overview.html}{MINT Platform} during a year-long competition with MITRE to provide DARPA the ``Airflow for science''
    \end{itemize}}

  \subsection{\textsc{Teaching}}
    \cventry{Jan.~2018 -- Dec.~2020}{Co-instructor}{The University of Texas at Austin}{Austin, TX}{}{
    \begin{itemize}
        \item Helped design and teach graduate-level courses: \emph{Machine Learning for the Geosciences} for Petrobras engineers, \emph{Intelligent Systems for the Geosciences}, and \emph{Scientific Computation}, including C++, CUDA, and HPC.
    \end{itemize}}

  \section{\textsc{Select Publications}}
    \cvitem{2021}{Gil, Y., et al. (incl. \textbf{D. Hardesty Lewis}), ``Artificial Intelligence for Modeling Complex Systems: Taming the Complexity of Expert Models to Improve Decision Making''. \emph{ACM TiiS}, 2021. DOI: \href{https://doi.org/10.1145/3453172}{10.1145/3453172}}
    \cvitem{2019}{Garijo, D., et al. (incl. \textbf{D. Hardesty Lewis}), ``An intelligent interface for integrating climate, hydrology, agriculture, and socioeconomic models''. \emph{ACM IUI '19}, 2019. DOI: \href{https://doi.org/10.1145/3308557.3308711}{10.1145/3308557.3308711}}

  \section{\textsc{Education}}
    \cventry{expected 2026}{M.S., Urban Planning}{Columbia University}{New York City,~NY}{}{}
    \cventry{2017}{B.S., Pure Mathematics}{University of Texas}{Austin,~TX}{}{}
    \cventry{2017}{Certificate, Scientific~Computation}{University of Texas}{Austin,~TX}{}{}

  \section{\textsc{Key Projects}}

    \cvitem{oppsAlert}{Developed the first property- and owner-level ML model of NIMBYism, using 15+ years of 15,000+ points of individual homeowner protest data provided by the City of Austin.}


  \section{\textsc{Activities}}
    \cvdoubleitem{Sailing}{Volunteer Atlantic SAR of the \href{https://en.wikipedia.org/wiki/Cheeki_Rafiki}{Cheeki Rafiki}}{Bikepacking}{US West \& East Coasts}

  \section{\textsc{Languages}}
    \cvdoubleitem{Spanish}{Native Proficiency}{French}{Professional Working Proficiency}

\end{document}
